\documentclass[11pt,a4paper]{article}

% ── Packages ──
\usepackage[UTF8, scheme=plain]{ctex}
\usepackage[T1]{fontenc}
\usepackage{amsmath,amssymb,amsthm}
\usepackage{mathtools}
\usepackage{bm}
\usepackage{graphicx}
\usepackage{booktabs}
\usepackage[colorlinks=true,citecolor=blue,linkcolor=blue]{hyperref}
\usepackage{geometry}
\geometry{margin=1.0in}
\usepackage[numbers,sort&compress]{natbib}
\usepackage{algorithm}
\usepackage{algpseudocode}
\usepackage{multirow}
\usepackage{xcolor}
\usepackage{enumitem}

% ── SCX Theorem / Rigor Environments ──
\newtheorem{theorem}{Theorem}[section]
\newtheorem{lemma}[theorem]{Lemma}
\newtheorem{corollary}[theorem]{Corollary}
\newtheorem{proposition}[theorem]{Proposition}
\theoremstyle{definition}
\newtheorem{definition}[theorem]{Definition}
\newtheorem{assumption}[theorem]{Assumption}
\theoremstyle{remark}
\newtheorem{remark}[theorem]{Remark}

% ── SCX Rigor Markers ──
\newcommand{\rigorFull}{\textcolor{green!60!black}{[rigor: full]}}
\newcommand{\rigorPartial}{\textcolor{orange!80!black}{[rigor: partial]}}
\newcommand{\rigorSketch}{\textcolor{red!60!black}{[rigor: sketch]}}
\newcommand{\assumptionTag}[1]{\textsf{\textbf{[A#1]}}}
\newcommand{\limitationTag}[1]{\textsf{\textbf{[L#1]}}}

% ── SCX Framework Macros ──
\newcommand{\SCX}{\textsc{SCX}}
\newcommand{\Situs}{\textsc{Situs}}
\newcommand{\Yajie}{\textsc{Yajie}}
\newcommand{\Cercis}{\textsc{Cercis}}
\newcommand{\Spring}{\textsc{Spring}}

% ── Math Operators ──
\DeclareMathOperator*{\argmin}{arg\,min}
\DeclareMathOperator*{\argmax}{arg\,max}
\DeclareMathOperator{\TV}{TV}
\DeclareMathOperator{\KL}{KL}
\DeclareMathOperator{\Var}{Var}
\DeclareMathOperator{\Cov}{Cov}
\DeclareMathOperator{\Corr}{Corr}
\DeclareMathOperator{\diag}{diag}
\DeclareMathOperator{\Bias}{Bias}
\DeclareMathOperator{\MSE}{MSE}
\DeclareMathOperator{\RMSE}{RMSE}
\DeclareMathOperator{\CRPS}{CRPS}
\newcommand{\R}{\mathbb{R}}
\newcommand{\N}{\mathbb{N}}
\newcommand{\E}{\mathbb{E}}
\newcommand{\Pbb}{\mathbb{P}}
\newcommand{\indicator}{\mathbf{1}}
\newcommand{\ind}[1]{\mathbf{1}[#1]}
\newcommand{\norm}[1]{\left\lVert#1\right\rVert}
\newcommand{\abs}[1]{\left|#1\right|}
\newcommand{\cD}{\mathcal{D}}
\newcommand{\cF}{\mathcal{F}}
\newcommand{\cG}{\mathcal{G}}
\newcommand{\cH}{\mathcal{H}}
\newcommand{\cK}{\mathcal{K}}
\newcommand{\cM}{\mathcal{M}}
\newcommand{\cN}{\mathcal{N}}
\newcommand{\cP}{\mathcal{P}}
\newcommand{\cS}{\mathcal{S}}
\newcommand{\cX}{\mathcal{X}}
\newcommand{\cY}{\mathcal{Y}}
\newcommand{\cZ}{\mathcal{Z}}

% ── Economics Domain Macros ──
\newcommand{\GDP}{\mathrm{GDP}}
\newcommand{\CPI}{\mathrm{CPI}}
\newcommand{\UNEMP}{\mathrm{UNEMP}}
\newcommand{\FFR}{\mathrm{FFR}}
\newcommand{\regime}{\boldsymbol{\rho}}
\newcommand{\statevec}{\mathbf{s}}
\newcommand{\obsvec}{\mathbf{y}}
\newcommand{\forecast}{\hat{\mathbf{y}}}
\newcommand{\shock}{\boldsymbol{\xi}}

% ── Existing SCX Theorems Reference ──
\newcommand{\ThmSCXNoise}{Theorem~S1}
\newcommand{\ThmSCXHonest}{Theorem~S3}

% ── Title ──
\title{\textbf{SCX-Audited Macroeconomic Forecasting:}\\
\large Multi-Model Consensus under Structural Breaks}

\author{SCX Research Group}
\date{Version 1.0 --- June 2026}

\begin{document}
\maketitle

\begin{abstract}
Macroeconomic forecasting faces a fundamental challenge: the economy is a non-stationary, regime-shifting system where the data-generating process itself changes discontinuously at structural breaks (financial crises, pandemics, policy regime shifts). Four families of forecasting models---Dynamic Stochastic General Equilibrium (DSGE), Vector Autoregressions (VAR), agent-based models (ABM), and machine learning (ML)---make different structural assumptions about economic dynamics, yet no formal framework certifies when their consensus or disagreement signals genuine structural change versus model failure. We apply the \SCX{} (Structured Causal eXamination) auditing framework to macroeconomic forecasting, providing: (i) a multi-model consensus detection theorem bounding the probability of missing systematic forecasting errors as $\exp(-2 M^{\mathrm{eff}} \Delta^2 / B^2)$ for $M$ forecasting experts; (ii) an unidentifiability theorem proving that model misspecification and genuine structural breaks (e.g., COVID-19, the 2008 financial crisis) are observationally equivalent without declared assumptions, with the Lucas Critique~\cite{Lucas1976} emerging as the special case where the structural break is a policy regime change; (iii) a \Cercis{} score $S = Q + \eta N$ where $Q$ combines RMSE accuracy and density calibration (via the Continuous Ranked Probability Score) and $N$ quantifies regime novelty relative to the historical training distribution; and (iv) a multi-model \Yajie{} consensus mechanism aggregating DSGE, VAR, ABM, and ML forecasts with explicit assumption tracking. All theorems are proved with explicit assumptions declared, and an experimental protocol is specified using FRED data with COVID-19 and 2008 financial crisis as structural break test cases. The framework establishes that macroeconomic forecast auditing requires---not merely benefits from---multi-model diversity, and that the impossibility of distinguishing model failure from structural change (without assumptions) is a feature that forces forecasters to declare their structural priors explicitly.

\medskip
\noindent\textbf{Keywords:} macroeconomic forecasting 宏观经济预测, structural breaks 结构突变, Lucas Critique 卢卡斯批判, SCX auditing, multi-model consensus 多模型共识, \Cercis{} scoring, \Yajie{} consensus, DSGE, VAR, agent-based models, machine learning, density calibration, FRED data
\end{abstract}

% ======================================================================
\section{Introduction 引言 宏观经济预测与结构突变}
% ======================================================================

Macroeconomic forecasting is the enterprise of predicting aggregate economic variables---Gross Domestic Product (GDP) growth, inflation (Consumer Price Index, CPI), unemployment rate, and the Federal Funds Rate (FFR)---over horizons ranging from one quarter to several years. The stakes are enormous: central bank interest rate decisions, fiscal stimulus packages, and institutional investment allocations all depend on these forecasts. Yet the track record of macroeconomic forecasting is sobering. Consensus forecasts systematically missed the 2008 Global Financial Crisis (GFC), the depth of the COVID-19 recession, and the persistence of post-pandemic inflation~\cite{Wieland2020, Ng2021}. The root cause is not any single forecasting model's failure, but a deeper structural problem: \textbf{the economy changes regime in ways that are not anticipated by any model trained on historical data}.

\subsection{The Macroeconomic Forecasting Landscape 宏观经济预测格局}

Four distinct modeling paradigms dominate macroeconomic forecasting, each embodying different---often contradictory---structural assumptions about how the economy operates:

\begin{enumerate}[label=(\roman*)]
    \item \textbf{Dynamic Stochastic General Equilibrium (DSGE) models 动态随机一般均衡.} These are microfounded models in which optimizing households and firms interact in markets that clear under rational expectations~\cite{SmetsWouters2007, ChristianoEichenbaumEvans2005}. DSGE models impose strong theoretical structure: parameters like the intertemporal elasticity of substitution, the Calvo price stickiness parameter, and the Taylor rule coefficients are estimated or calibrated from long-run economic relationships. Their forecasts are disciplined by economic theory but fragile to structural breaks that invalidate the underlying behavioral assumptions.
    
    \item \textbf{Vector Autoregressions (VAR) 向量自回归.} VAR models---and their Bayesian (BVAR) and time-varying parameter (TVP-VAR) extensions~\cite{Sims1980, Primiceri2005}---treat the economy as a linear (or smoothly time-varying) stochastic process. They make minimal structural assumptions, instead capturing statistical regularities in the historical comovement of macroeconomic variables. VARs are flexible but extrapolate past correlations, making them vulnerable to regime changes where historical comovement patterns break down.

    \item \textbf{Agent-Based Models (ABM) 基于主体的模型.} ABMs simulate heterogeneous agents (firms, households, banks) interacting through decentralized markets with bounded rationality~\cite{HaldaneTurrell2018, Dosi2010}. They can generate emergent phenomena---financial bubbles, sudden recessions, nonlinear propagation of shocks---that aggregate models miss. However, ABMs have vast parameter spaces, making calibration difficult and forecast uncertainty large.

    \item \textbf{Machine Learning (ML) models 机器学习模型.} Neural networks, gradient boosting, random forests, and recurrent architectures have been applied to macroeconomic forecasting with notable success~\cite{Coulombe2020, Medeiros2021}. ML models capture nonlinear interactions and high-dimensional predictors but lack structural interpretability. Their out-of-distribution behavior---precisely when the economy undergoes a structural break---is unguaranteed.
\end{enumerate}

Each of these model families constitutes an ``expert'' in the \SCX{} sense: a distinct forecasting function $f_m: \cX \to \cY$ mapping economic states to predictions, with distinct inductive biases, training procedures, and structural assumptions.

\subsection{The Lucas Critique as a Canonical Unidentifiability Problem 卢卡斯批判作为典型的不可辨识性问题}

The Lucas Critique~\cite{Lucas1976} is the foundational statement of the structural break problem in macroeconomics. Lucas argued that econometric models estimated on historical data cannot be used to evaluate alternative policy regimes because the parameters of those models---being functions of agents' decision rules---will change when the policy regime changes. Formally: a forecasting model $f_\theta$ with parameters $\theta$ estimated under policy regime $\pi_1$ will produce biased forecasts under policy regime $\pi_2 \neq \pi_1$, because the true parameter $\theta^*(\pi)$ is a function of the policy regime that the model treats as constant.

The Lucas Critique, we argue, is a special case of a more general unidentifiability principle. When a forecasting model's predictions deviate from realized outcomes, the deviation can arise from:
\begin{enumerate}[label=(\roman*)]
    \item \textbf{Model misspecification} (the model is wrong even for the historical regime);
    \item \textbf{Structural break} (the data-generating process changed, and no model trained on pre-break data could have predicted the post-break outcomes);
    \item \textbf{Measurement error} (the realized data are noisy estimates of the true economic state);
    \item \textbf{Policy regime change} (the Lucas Critique case: agents' decision rules changed because policy changed).
\end{enumerate}
Without declared assumptions about which of these sources is bounded, post-hoc error attribution is logically underdetermined. \textbf{The Lucas Critique is the special case where the structural break is explicitly identified as a policy regime change}, but the broader principle applies to any structural break.

\subsection{Our Contribution}

We apply the \SCX{} auditing framework to macroeconomic forecasting. The core thesis is:

\begin{center}
\emph{Multi-model macroeconomic forecasting is not merely a best practice for robustness.\\
It is a logical necessity for detecting when individual models fail---\\
and for honestly acknowledging that model failure and structural change are\\
fundamentally unidentifiable without declared assumptions.}
\end{center}

The paper's contributions are:

\begin{enumerate}[label=(\roman*)]
    \item \textbf{Formalization} (\S2): Economic regime as state, forecasting model as expert, structural break as regime transition.
    \item \textbf{Theorem 1 --- Multi-Model Forecast Error Detection} (\S3): A bound on the probability of missing systematic forecasting errors when $M$ models are deployed simultaneously.
    \item \textbf{Theorem 2 --- Model Misspecification vs.\ Structural Break Unidentifiability} (\S4): Formal proof that model error and structural change are observationally equivalent without declared assumptions, with Lucas Critique as a corollary.
    \item \textbf{\Cercis{} Score} (\S5): Quality score combining forecast accuracy (RMSE, CRPS density calibration) with regime novelty for ranking forecasting models.
    \item \textbf{Multi-Model \Yajie{} Consensus} (\S6): Game-theoretic consensus mechanism aggregating DSGE, VAR, ABM, and ML forecasts with explicit assumption tracking.
    \item \textbf{Experimental Protocol} (\S7): Benchmark specification using FRED data with COVID-19 and 2008 GFC as structural break test cases.
    \item \textbf{Discussion} (\S8): Implications for central bank forecasting, honest uncertainty communication, and the limits of model-based prediction.
\end{enumerate}

\assumptionTag{0} \textbf{(Motivating Observation)}: No existing macroeconomic forecasting framework provides a formal guarantee on the probability of missing a systematic forecasting error, nor a formal analysis of when model disagreement signals structural change versus model failure. The \SCX{} framework fills this gap.

\limitationTag{0} The framework is descriptive and diagnostic; it does not propose new forecasting models. It provides the quality audit infrastructure that any forecasting model---existing or future---can be evaluated against.

% ======================================================================
\section{Formalization: Macroeconomic Forecasting as State-Conditioned Expert Prediction 宏观经济预测作为状态条件专家预测的形式化}
\label{sec:formalization}
% ======================================================================

\subsection{The Economic State 经济状态}

\begin{definition}[Economic State Space 经济状态空间]
\label{def:state}
The economic state at time $t$ is a vector:
\begin{equation}
\statevec_t = (\mathbf{y}_{t-h:t}, \mathbf{x}_{t-h:t}, \regime_t) \in \cS = \R^{d_y \cdot h} \times \R^{d_x \cdot h} \times \mathcal{R},
\label{eq:state}
\end{equation}
where:
\begin{itemize}[nosep]
    \item $\mathbf{y}_{t-h:t} \in \R^{d_y \cdot h}$: the history of $d_y$ macroeconomic variables over the past $h$ periods (typically $h = 8$ quarters for a 2-year lookback);
    \item $\mathbf{x}_{t-h:t} \in \R^{d_x \cdot h}$: the history of $d_x$ exogenous or conditioning variables (global GDP, commodity prices, financial conditions indices);
    \item $\regime_t \in \mathcal{R}$: the current \textbf{economic regime}---a latent categorical variable capturing the qualitative state of the economy. We define $\mathcal{R} = \{r_{\mathrm{expansion}}, r_{\mathrm{recession}}, r_{\mathrm{crisis}}, r_{\mathrm{recovery}}, r_{\mathrm{stagflation}}, r_{\mathrm{pandemic}}\}$.
\end{itemize}
\end{definition}

\begin{definition}[Structural Break 结构突变]
\label{def:structural_break}
A \textbf{structural break} at time $T$ is a discontinuous change in the data-generating process (DGP). Formally, there exists $T$ such that for all $t \geq T$, the conditional distribution of $\mathbf{y}_{t+1:t+H}$ given $\statevec_t$ differs from its pre-break distribution:
\begin{equation}
\Pbb(\mathbf{y}_{t+1:t+H} \mid \statevec_t) \neq \Pbb^{\mathrm{pre}}(\mathbf{y}_{t+1:t+H} \mid \statevec_t), \quad \forall t \geq T.
\label{eq:structural_break}
\end{equation}
The regime variable satisfies $\regime_T \neq \regime_{T-1}$ at the break point.
\end{definition}

\assumptionTag{1} \textbf{(Regime Identifiability Ex Post)}: The regime $\regime_t$ is identifiable after a delay of $D$ periods: at time $t + D$, an economic historian can assign $\regime_t \in \mathcal{R}$ with high confidence. For example, the NBER Business Cycle Dating Committee identifies recession start dates with a typical lag of 6--18 months. This assumption grounds the regime concept in observable practice.

\subsection{Forecasting Models as Experts 预测模型作为专家}

\begin{definition}[Forecasting Expert 预测专家]
\label{def:expert}
A forecasting expert is a function $f_m: \cS \to \cY$ mapping the current economic state to a forecast of the next $H$ periods:
\begin{equation}
f_m(\statevec_t) = (\hat{y}_{t+1}^{(m)}, \hat{y}_{t+2}^{(m)}, \ldots, \hat{y}_{t+H}^{(m)}) \in \cY \subset \R^{d_y \cdot H},
\label{eq:forecast}
\end{equation}
where $\hat{y}_{t+h}^{(m)} = (\widehat{\GDP}_{t+h}, \widehat{\CPI}_{t+h}, \widehat{\UNEMP}_{t+h}, \widehat{\FFR}_{t+h})^{(m)}$ is the $m$-th model's $h$-step-ahead forecast of the four key macroeconomic variables. Probabilistic experts additionally output a predictive distribution $\hat{P}_m(\mathbf{y}_{t+1:t+H} \mid \statevec_t)$.
\end{definition}

We consider $M$ experts partitioned into four families:
\begin{equation}
\cF = \cF_{\mathrm{DSGE}} \cup \cF_{\mathrm{VAR}} \cup \cF_{\mathrm{ABM}} \cup \cF_{\mathrm{ML}},
\label{eq:families}
\end{equation}
where $|\cF_{\mathrm{DSGE}}| = M_{\mathrm{DSGE}}$, etc., and $M = \sum_k M_k$.

\begin{definition}[Forecast Error 预测误差]
\label{def:forecast_error}
For expert $m$ at lead time $h$, the forecast error is:
\begin{equation}
\varepsilon_{t,h}^{(m)} = \norm{\hat{y}_{t+h}^{(m)} - y_{t+h}}_{\Sigma},
\label{eq:forecast_error}
\end{equation}
where $\norm{\mathbf{v}}_{\Sigma} = \sqrt{\mathbf{v}^T \Sigma^{-1} \mathbf{v}}$ is the Mahalanobis norm with covariance $\Sigma = \Cov(\mathbf{y})$ estimated from historical data. This normalization ensures that errors in GDP (units: \% change) and inflation (units: percentage points) are comparable.
\end{definition}

\assumptionTag{2} \textbf{(Bounded Output)}: Macroeconomic variables are bounded in practice: GDP growth $\in [-20\%, +20\%]$ at quarterly frequency, CPI inflation $\in [-5\%, +25\%]$, unemployment $\in [2\%, 25\%]$, FFR $\in [0\%, 20\%]$. This implies a compact output space $\cY \subset \R^{d_y \cdot H}$ with diameter $B < \infty$.

\assumptionTag{3} \textbf{(Expert Diversity)}: Experts from different families make genuinely different structural assumptions. Formally, for any structural break $\regime_t \neq \regime_{t-1}$, the correlation of forecast errors across families satisfies $\Corr(\varepsilon^{(m)}, \varepsilon^{(m')}) \leq \rho_{\max} < 1$ for $m \in \cF_k, m' \in \cF_{k'}, k \neq k'$. Within-family correlation may be high due to shared structural assumptions.

\subsection{The Lucas Critique in This Framework 卢卡斯批判在此框架中的表述}

\begin{proposition}[Lucas Critique as a Regime-Dependent Parameter 卢卡斯批判作为制度依赖参数]
\label{prop:lucas_formal}
Let a DSGE model $f_{\mathrm{DSGE}}$ have structural parameters $\boldsymbol{\theta} \in \Theta$ estimated under policy regime $\pi$. Under a new policy regime $\pi'$, the true structural parameters shift to $\boldsymbol{\theta}^*(\pi') \neq \boldsymbol{\theta}^*(\pi)$. The DSGE forecast error decomposes as:
\begin{equation}
\varepsilon_{\mathrm{DSGE}} = \underbrace{\varepsilon_{\mathrm{param}}}_{\text{parameter shift}} + \underbrace{\varepsilon_{\mathrm{spec}}}_{\text{model misspecification}} + \underbrace{\varepsilon_{\mathrm{shock}}}_{\text{unforecastable shock}},
\label{eq:lucas_decomp}
\end{equation}
where $\varepsilon_{\mathrm{param}} = \norm{f_{\boldsymbol{\theta}^*(\pi)}(\statevec) - f_{\boldsymbol{\theta}^*(\pi')}(\statevec)}$ is the error attributable to the policy regime change. The Lucas Critique states that $\varepsilon_{\mathrm{param}} > 0$ when $\pi' \neq \pi$, and that this component is identifiable only if the policy regime change is observed and the mapping $\boldsymbol{\theta}^*(\cdot)$ is known.
\end{proposition}

\begin{proof}\rigorPartial
This is a restatement of Lucas's (1976) argument in our notation. The key insight is that $\boldsymbol{\theta}^*(\pi) \neq \boldsymbol{\theta}^*(\pi')$ follows from agents re-optimizing under the new policy rule, but the new parameters are not estimable from pre-change data. $\square$
\end{proof}

% ======================================================================
\section{Theorem 1: Multi-Model Forecast Error Detection 多模型预测误差检测定理}
\label{sec:noise_detection}
% ======================================================================

We now prove the core detection theorem: when $M$ forecasting models are deployed simultaneously, the probability of missing a systematic forecasting error decays exponentially in $M$.

\subsection{Setup}

Consider $M$ forecasting experts $\{f_1, \ldots, f_M\}$ evaluated at time $t$ on the same economic state $\statevec_t$. Let $y_{t+h}^*$ denote the realized value of the macroeconomic vector at time $t+h$. An expert $m$ has a \textbf{systematic error} of magnitude $\Delta > 0$ at horizon $h$ if:
\begin{equation}
\norm{\hat{y}_{t+h}^{(m)} - y_{t+h}^*}_{\Sigma} > \Delta.
\label{eq:systematic_error}
\end{equation}

\begin{definition}[Pairwise Discrepancy]
\label{def:discrepancy}
The pairwise discrepancy between experts $i$ and $j$ at time $t$, horizon $h$ is:
\begin{equation}
d_{ij}^{(t,h)} = \norm{\hat{y}_{t+h}^{(i)} - \hat{y}_{t+h}^{(j)}}_{\Sigma}.
\label{eq:discrepancy}
\end{equation}
\end{definition}

\assumptionTag{4} \textbf{(Detectability Condition 可检测条件)}: If an expert $i$ has systematic error exceeding $\Delta$ at state $\statevec_t$, then for at least a fraction $\gamma \in (0, 1]$ of other experts $j$ from \textit{different model families}, we have $d_{ij}^{(t,h)} > \Delta/2$.

\textit{Justification:} If expert $i$ errs by more than $\Delta$ and most other experts are within $\Delta/2$ of the truth, then by the reverse triangle inequality: $d_{ij} = \norm{\hat{y}^{(i)} - \hat{y}^{(j)}} \geq \norm{\hat{y}^{(i)} - y^*} - \norm{\hat{y}^{(j)} - y^*} > \Delta - \Delta/2 = \Delta/2$. The cross-family requirement ensures that this holds even when within-family experts share the same error (due to shared structural assumptions). This is the macroeconomic analog of the CFD detectability condition from the SCX framework and formalizes the intuition that a ``bad'' forecast looks different from most ``good'' forecasts.

\subsection{Theorem Statement}

\begin{theorem}[Multi-Model Forecast Error Detection 多模型预测误差检测]
\label{thm:macro_noise}
\rigorFull
Let $M \geq 3$ forecasting experts from $K$ distinct model families be applied at time $t$. Let $M_{\mathrm{cross}}^{(i)}$ be the number of experts from families other than expert $i$'s own family. Assume \assumptionTag{4} holds with detectability fraction $\gamma$. Define the consensus indicator for expert $i$:
\begin{equation}
C_i^{(t)} = \indicator\!\left( \frac{1}{M_{\mathrm{cross}}^{(i)}} \sum_{j \in \cF \setminus \cF_{(i)}} \indicator(d_{ij}^{(t,h)} \leq \Delta/2) \geq 1 - \theta \right),
\label{eq:consensus_indicator}
\end{equation}
where $\theta \in (0, 1-\gamma]$ is a consensus threshold and $\cF_{(i)}$ is expert $i$'s model family.

Let $E_i$ be the event that expert $i$ has systematic error exceeding $\Delta$ but is not detected (i.e., $C_i^{(t)} = 1$ despite the error). Then:
\begin{equation}
\Pbb(E_i) \leq \exp\!\left( -2 M_{\mathrm{cross}}^{(i)} (\gamma - \theta)^2 \right).
\label{eq:macro_noise_bound}
\end{equation}
\end{theorem}

\begin{proof}
Fix a time $t$ where expert $i$ has systematic error $> \Delta$. By \assumptionTag{4}, for at least $\gamma M_{\mathrm{cross}}^{(i)}$ cross-family experts $j$, we have $d_{ij}^{(t,h)} > \Delta/2$.

Define the binary variable $X_j = \indicator(d_{ij}^{(t,h)} \leq \Delta/2)$ for each cross-family expert $j$. Under the error condition, at most $(1-\gamma) M_{\mathrm{cross}}^{(i)}$ of these satisfy $X_j = 1$. Let $\bar{X} = \frac{1}{M_{\mathrm{cross}}^{(i)}} \sum_{j} X_j$ denote the fraction of cross-family experts whose predictions are ``close'' (within $\Delta/2$) to expert $i$'s (erroneous) prediction.

For a false negative (missing the error), we need $\bar{X} \geq 1 - \theta$, meaning at least a fraction $1-\theta$ of cross-family experts agree with the erroneous expert. Since under the error condition $\E[\bar{X}] \leq 1 - \gamma$, this requires $\bar{X} - \E[\bar{X}] \geq (1 - \theta) - (1 - \gamma) = \gamma - \theta$.

The $X_j$ are bounded in $[0, 1]$ and, for cross-family experts, are independent conditional on the state (since different model families use independent data representations, training procedures, and structural assumptions). Applying Hoeffding's inequality~\cite{Hoeffding1963}:
\begin{equation}
\Pbb(\bar{X} - \E[\bar{X}] \geq \gamma - \theta) \leq \exp\!\left( -2 M_{\mathrm{cross}}^{(i)} (\gamma - \theta)^2 \right).
\label{eq:hoeffding}
\end{equation}

Substituting $\bar{X} \geq 1-\theta$ as the false negative condition yields the bound. $\square$
\end{proof}

\begin{corollary}[Required Expert Multiplicity 所需专家多重性]
\label{cor:required_M}
To achieve detection confidence $1 - \alpha$ (i.e., $\Pbb(E_i) \leq \alpha$) with $\gamma = 0.5$ (errors detectable by at least half of cross-family experts) and $\theta = 0.2$ (allow up to 20\% false agreement), the required cross-family multiplicity is:
\begin{equation}
M_{\mathrm{cross}}^{(i)} \geq \frac{\log(1/\alpha)}{2(\gamma - \theta)^2} = \frac{\log(1/\alpha)}{2 \cdot 0.3^2} \approx 5.56 \cdot \log(1/\alpha).
\label{eq:required_M}
\end{equation}
For $\alpha = 0.01$ (99\% confidence), $M_{\mathrm{cross}}^{(i)} \geq 26$. With $K = 4$ model families, this requires approximately $M_k \approx 7$ experts per family, yielding $M = 28$ total experts.
\end{corollary}

\begin{corollary}[Joint Detection Across Horizons 跨预测期的联合检测]
\label{cor:joint_detection}
For $H$ forecast horizons evaluated at time $t$, the probability of missing a systematic error at \textit{any} horizon is bounded by the union bound:
\begin{equation}
\Pbb(\text{miss at any } h) \leq H \cdot \exp\!\left( -2 M_{\mathrm{cross}}^{(i)} (\gamma - \theta)^2 \right).
\label{eq:joint_bound}
\end{equation}
For $H = 8$ (quarters, 2-year forecast horizon), $M_{\mathrm{cross}}^{(i)} = 20$, $\gamma = 0.5$, $\theta = 0.2$: $\Pbb(\text{miss}) \leq 8 \cdot \exp(-2 \cdot 20 \cdot 0.09) = 8 \cdot \exp(-3.6) \approx 8 \cdot 0.0273 \approx 0.219$, which provides a non-trivial (though not extremely tight) bound. Tighter bounds require larger $M_{\mathrm{cross}}^{(i)}$ or a Slepian-type correction for dependent horizons.
\end{corollary}

\limitationTag{1} The Hoeffding bound assumes independent expert errors across model families. Experts within the same family (e.g., two BVAR specifications with different priors) may exhibit correlated errors. The cross-family restriction $M_{\mathrm{cross}}^{(i)}$ partially addresses this, but residual correlation can be handled via the effective sample size $M_{\mathrm{eff}} = M_{\mathrm{cross}}^{(i)} / (1 + (M_{\mathrm{cross}}^{(i)}-1) \bar{\rho})$ where $\bar{\rho}$ is the average pairwise error correlation across families~\cite{Liang1986}.

\begin{remark}[Why Cross-Family Diversity Matters 为什么跨模型族多样性重要]
\label{rem:diversity}
Theorem~\ref{thm:macro_noise} reveals a structural requirement for macroeconomic forecast auditing: \textbf{the detection guarantee depends on cross-family expert diversity, not merely on the total number of experts.} Deploying 50 VAR variants provides weaker detection than deploying 5 DSGE + 5 VAR + 5 ABM + 5 ML experts because within-family experts respond similarly to structural breaks. This is the formal justification for the multi-paradigm approach advocated by the SCX framework and by central bank forecast combination practices~\cite{AiolfiTimmermann2006}.
\end{remark}

% ======================================================================
\section{Theorem 2: Model Misspecification vs.\ Structural Break Unidentifiability 模型误设与结构突变的不可辨识性}
\label{sec:unidentifiability}
% ======================================================================

We now prove the paper's central theoretical result: when a forecasting model's predictions deviate from realized outcomes, the deviation source---model misspecification or genuine structural break---is fundamentally unidentifiable without declared assumptions. This generalizes the Lucas Critique from policy regime changes to all structural breaks.

\subsection{The Error Source Decomposition 误差来源分解}

Consider time $T$ when the economy undergoes a structural break (e.g., Q3 2008 for the GFC, Q1 2020 for COVID-19). For any forecasting model $f_m$ trained on pre-break data $\cD_{\mathrm{pre}} = \{(\statevec_t, \mathbf{y}_{t+1:t+H})\}_{t=1}^{T-1}$, the post-break forecast error at horizon $h$ decomposes as:

\begin{equation}
\varepsilon_{T,h}^{(m)} = \underbrace{\varepsilon_{\mathrm{spec}}^{(m)}}_{\text{model misspecification}} + \underbrace{\varepsilon_{\mathrm{break}}}_{\text{structural break}} + \underbrace{\varepsilon_{\mathrm{noise}}}_{\text{irreducible shock}},
\label{eq:error_decomp_macro}
\end{equation}

where:
\begin{itemize}[nosep]
    \item $\varepsilon_{\mathrm{spec}}^{(m)} = \norm{\E[f_m(\statevec_T) \mid \text{no break}] - \E[\mathbf{y}_{T+h} \mid \statevec_T, \text{no break}]}_{\Sigma}$: the model's error under the counterfactual of no structural break;
    \item $\varepsilon_{\mathrm{break}} = \norm{\E[\mathbf{y}_{T+h} \mid \statevec_T, \text{break}] - \E[\mathbf{y}_{T+h} \mid \statevec_T, \text{no break}]}_{\Sigma}$: the shift in the conditional mean due to the structural break;
    \item $\varepsilon_{\mathrm{noise}} = \norm{\mathbf{y}_{T+h} - \E[\mathbf{y}_{T+h} \mid \statevec_T, \text{break}]}_{\Sigma}$: the irreducible forecast error (the innovation).
\end{itemize}

\begin{assumption}[A5: Irreducible Noise is Unforecastable]
\label{ass:A5}
The innovation $\varepsilon_{\mathrm{noise}}$ is mean-zero and independent of the model choice: $\E[\varepsilon_{\mathrm{noise}} \mid f_m] = 0$ for all $m$. This is the standard rational expectations assumption that forecast errors are unforecastable~\cite{MincerZarnowitz1969}.
\end{assumption}

\subsection{Observational Equivalence Construction 观测等价构造}

\begin{theorem}[Model Misspecification vs.\ Structural Break Unidentifiability 模型误设与结构突变不可辨识性]
\label{thm:macro_unident}
\rigorFull
Let $\cF = \{f_1, \ldots, f_M\}$ be a set of forecasting models trained on pre-break data $\cD_{\mathrm{pre}}$. At a post-break time $T$, let the observed forecast errors be $\{\varepsilon_{T,h}^{(m)}\}_{m=1}^M$. Under \assumptionTag{5}, there exist at least two distinct decompositions of the error vector $\boldsymbol{\varepsilon}_T = (\varepsilon_{T,h}^{(1)}, \ldots, \varepsilon_{T,h}^{(M)})$ into model misspecification and structural break components that are \textbf{observationally equivalent}---they produce identical observed errors but attribute them to different causes. Consequently, without declared assumptions bounding either $\varepsilon_{\mathrm{spec}}^{(m)}$ or $\varepsilon_{\mathrm{break}}$, the attribution is unidentifiable.
\end{theorem}

\begin{proof}
We construct two observationally equivalent worlds.

\textbf{World A (Model Blame 模型归因):} Assume there was \textit{no structural break}: $\varepsilon_{\mathrm{break}} = 0$. Then the entire observed error for each model $m$ is attributed to model misspecification:
\begin{equation}
\varepsilon_{\mathrm{spec}}^{(m,A)} = \varepsilon_{T,h}^{(m)}, \quad \varepsilon_{\mathrm{break}}^{(A)} = 0, \quad \forall m.
\label{eq:world_A}
\end{equation}
Interpretation: ``Our models are simply bad at forecasting---they would have failed even without the pandemic/crisis.''

\textbf{World B (Break Blame 结构突变归因):} Assume all models are \textit{correctly specified for the pre-break regime}: $\varepsilon_{\mathrm{spec}}^{(m)} = 0$ for all $m$ up to a family-specific residual. Then the observed error is attributed entirely to the structural break:
\begin{equation}
\varepsilon_{\mathrm{spec}}^{(m,B)} = 0, \quad \varepsilon_{\mathrm{break}}^{(B)} = \frac{1}{M} \sum_{m=1}^M \varepsilon_{T,h}^{(m)}, \quad \forall m.
\label{eq:world_B}
\end{equation}
with the within-family dispersion $\varepsilon_{T,h}^{(m)} - \varepsilon_{\mathrm{break}}^{(B)}$ attributed to irreducible noise $\varepsilon_{\mathrm{noise}}$. Interpretation: ``No model could have predicted this---the economy fundamentally changed.''

\textbf{Observational equivalence:} Both worlds produce identical total forecast errors:
\begin{equation}
\varepsilon_{T,h}^{(m)} = \varepsilon_{\mathrm{spec}}^{(m,A)} + 0 = 0 + \varepsilon_{\mathrm{break}}^{(B)} + (\varepsilon_{T,h}^{(m)} - \varepsilon_{\mathrm{break}}^{(B)}).
\label{eq:obs_equiv}
\end{equation}
The first equality is World A; the second is World B with the residual assigned to $\varepsilon_{\mathrm{noise}}$. Since $\E[\varepsilon_{\mathrm{noise}} \mid f_m] = 0$ by \assumptionTag{5}, the mean error is preserved across both attributions.

\textbf{Continuum of intermediate worlds:} For any $\alpha \in [0, 1]$, define:
\begin{equation}
\varepsilon_{\mathrm{spec}}^{(m,\alpha)} = \alpha \cdot \varepsilon_{T,h}^{(m)}, \quad \varepsilon_{\mathrm{break}}^{(\alpha)} = (1-\alpha) \cdot \frac{1}{M} \sum_{m=1}^M \varepsilon_{T,h}^{(m)},
\label{eq:continuum}
\end{equation}
with the residual again assigned to noise. Every $\alpha$ produces an observationally equivalent attribution. The set of valid attributions is a $(M+1)$-dimensional simplex intersected with the observed error constraints.

\textbf{Dimensionality argument:} The full attribution vector $(\varepsilon_{\mathrm{spec}}^{(1)}, \ldots, \varepsilon_{\mathrm{spec}}^{(M)}, \varepsilon_{\mathrm{break}})$ has $M+1$ degrees of freedom. The observed data provide only $M$ constraints (one error per model). Without additional assumptions (providing at least one more constraint), the system is underdetermined with at least one degree of freedom, yielding an infinite family of valid attributions. $\square$
\end{proof}

\subsection{The Lucas Critique as a Special Case 卢卡斯批判作为特例}

\begin{corollary}[Lucas Critique as Identified Structural Break 卢卡斯批判作为可辨识的结构突变]
\label{cor:lucas_special}
The Lucas Critique is the special case of Theorem~\ref{thm:macro_unident} where the structural break is \textbf{identified} by an additional declared assumption: the policy regime change $\pi \to \pi'$ is \textit{observed} (e.g., a central bank announces a new inflation target or a new monetary policy rule), and the mapping $\boldsymbol{\theta}^*(\pi)$ is \textit{declared known} up to estimation uncertainty. Under these two assumptions:
\begin{enumerate}[label=(\roman*)]
    \item The structural break component $\varepsilon_{\mathrm{break}}$ is partially identifiable as $\varepsilon_{\mathrm{param}} = \norm{f_{\boldsymbol{\theta}^*(\pi)}(\statevec) - f_{\boldsymbol{\theta}^*(\pi')}(\statevec)}_{\Sigma}$;
    \item The residual model misspecification is $\varepsilon_{\mathrm{spec}} = \varepsilon_{\mathrm{total}} - \varepsilon_{\mathrm{param}} - \varepsilon_{\mathrm{noise}}$.
\end{enumerate}
Without the declaration that $\pi \to \pi'$ occurred and that $\boldsymbol{\theta}^*(\cdot)$ is known, the Lucas Critique error is unidentifiable---precisely the content of Theorem~\ref{thm:macro_unident}.
\end{corollary}

\begin{proof}
The additional assumptions (observed policy change, known parameter mapping) provide the $(M+2)$-th constraint needed to resolve the underdetermined system. When $\pi \to \pi'$ is unobserved or the parameter mapping is unknown, we return to the Theorem~\ref{thm:macro_unident} setting with $M+1$ degrees of freedom and $M$ constraints. $\square$
\end{proof}

\begin{remark}[COVID-19: Identifiable or Unidentifiable? COVID-19：可辨识还是不可辨识？]
\label{rem:covid}
The COVID-19 pandemic of 2020 illustrates the theorem's practical force. Consider GDP forecasts for Q2 2020 made in Q1 2020:
\begin{itemize}
    \item \textbf{World A (model failure):} ``Our DSGE model failed because it doesn't model lockdowns or epidemiological dynamics. This is model misspecification.''
    \item \textbf{World B (structural break):} ``No macroeconomic model could have forecasted a 30\% annualized GDP decline because a pandemic-triggered economic shutdown is a genuinely unprecedented structural break outside the support of any training data.''
\end{itemize}
Both claims are observationally equivalent from the forecast error alone. Theorem~\ref{thm:macro_unident} proves that \textit{resolving this debate requires declared assumptions}, not more data or better models. The honest position is: ``Given our assumption that pandemics are exogenous structural breaks outside the model class, we attribute the Q2 2020 forecast error to a structural break. Under the alternative assumption that models should incorporate pandemic dynamics, the error is model misspecification.''
\end{remark}

\subsection{Implications for Macroeconomic Forecast Auditing 对宏观经济预测审计的启示}

\begin{corollary}[Assumption Mandate for Forecast Error Attribution 预测误差归因的假设声明要求]
\label{cor:assumption_mandate_macro}
Any claim attributing a macroeconomic forecast error to model misspecification, structural break, or measurement error \textbf{must} be accompanied by explicit, falsifiable assumptions about the other two components. The \SCX{} audit requires forecasters to declare:
\begin{enumerate}[label=(\roman*)]
    \item \textbf{Model scope assumption:} ``Our model class $\cF_k$ is assumed to correctly specify the conditional mean $\E[\mathbf{y}_{t+h} \mid \statevec_t]$ for regimes in the set $\mathcal{R}_{\mathrm{covered}} \subseteq \mathcal{R}$.''
    \item \textbf{Break detection assumption:} ``We declare a structural break when [criterion], and we assume the break magnitude $\varepsilon_{\mathrm{break}}$ is bounded above by $B_{\mathrm{break}}$.''
    \item \textbf{Measurement assumption:} ``We assume the realized data $\mathbf{y}_{t+h}$ measure the true economic state with error bounded by $\sigma_{\mathrm{meas}}$.''
\end{enumerate}
Without these declarations, post-hoc error attribution is logically underdetermined (by Theorem~\ref{thm:macro_unident}) and therefore scientifically invalid.
\end{corollary}

% ======================================================================
\section{\Cercis{} Score for Macroeconomic Forecasts 宏观经济预测的Cercis评分}
\label{sec:cercis}
% ======================================================================

The \Cercis{} scoring framework ranks forecasting models along two axes: quality $Q$ (accuracy on benchmarks) and novelty $N$ (coverage of economic regimes). We adapt it to macroeconomic forecasting.

\subsection{Quality Score $Q$: Forecast Accuracy and Density Calibration 预测准确度与密度校准}

\begin{definition}[Macroeconomic Quality Score 宏观经济质量评分]
\label{def:Q_macro}
For expert $f_m$ evaluated over a set of $T$ time periods, the quality score combines point forecast accuracy and probabilistic calibration:
\begin{equation}
Q(f_m) = \underbrace{Q_{\mathrm{RMSE}}(f_m)}_{\text{point accuracy}} \;+\; \lambda_{\mathrm{cal}} \cdot \underbrace{Q_{\mathrm{cal}}(f_m)}_{\text{density calibration}},
\label{eq:Q_macro}
\end{equation}
where $\lambda_{\mathrm{cal}} \geq 0$ weights the importance of probabilistic calibration relative to point accuracy.

\textbf{Point Accuracy (RMSE-based):}
\begin{equation}
Q_{\mathrm{RMSE}}(f_m) = 1 - \frac{1}{T} \sum_{t=1}^T \min\!\left(1, \frac{\RMSE_{t}^{(m)}}{\tau_{\mathrm{RMSE}}}\right),
\label{eq:Q_RMSE}
\end{equation}
where $\RMSE_{t}^{(m)} = \sqrt{\frac{1}{H d_y} \sum_{h=1}^H \norm{\hat{y}_{t+h}^{(m)} - y_{t+h}}_{\Sigma}^2}$ is the root mean squared error across $H$ horizons and $d_y$ variables, and $\tau_{\mathrm{RMSE}}$ is a tolerance threshold (e.g., $\tau_{\mathrm{RMSE}} = 2.0$, corresponding to errors twice the historical standard deviation).

\textbf{Density Calibration (CRPS-based):}
For probabilistic forecasts outputting a predictive distribution $\hat{P}_m(\mathbf{y}_{t+h} \mid \statevec_t)$, calibration is assessed via the Continuous Ranked Probability Score~\cite{GneitingRaftery2007}:
\begin{equation}
\CRPS(\hat{P}_m, y_{t+h}) = \int_{-\infty}^{\infty} (\hat{F}_m(z) - \indicator(z \geq y_{t+h}))^2 dz,
\label{eq:CRPS}
\end{equation}
where $\hat{F}_m$ is the predictive CDF. The calibration score normalizes CRPS against the climatological baseline (unconditional distribution):
\begin{equation}
Q_{\mathrm{cal}}(f_m) = 1 - \frac{\frac{1}{T H} \sum_{t=1}^T \sum_{h=1}^H \CRPS(\hat{P}_m^{(t,h)}, y_{t+h})}{\frac{1}{T H} \sum_{t=1}^T \sum_{h=1}^H \CRPS(\hat{P}_{\mathrm{clim}}, y_{t+h})}.
\label{eq:Q_cal}
\end{equation}
$Q_{\mathrm{cal}} \in (-\infty, 1]$: positive values indicate skill relative to climatology; zero means no skill; negative values indicate the forecast is worse than the unconditional distribution.
\end{definition}

\begin{definition}[Structural Break Penalty 结构突变惩罚]
\label{def:break_penalty}
To prevent models from achieving high $Q$ by avoiding difficult regimes, we compute a regime-conditioned quality:
\begin{equation}
Q_{\mathrm{regime}}(f_m, r) = Q(f_m \mid \regime_t = r),
\label{eq:Q_regime}
\end{equation}
and define the regime-weighted quality:
\begin{equation}
\tilde{Q}(f_m) = \sum_{r \in \mathcal{R}} w_r \cdot Q_{\mathrm{regime}}(f_m, r),
\label{eq:Q_weighted}
\end{equation}
where $w_r$ are regime importance weights. By default, $w_{\mathrm{crisis}} = w_{\mathrm{pandemic}} = 0.3$, $w_{\mathrm{recession}} = 0.2$, $w_{\mathrm{expansion}} = w_{\mathrm{recovery}} = 0.1$, reflecting the higher importance of forecasting accuracy during turbulent regimes.
\end{definition}

\subsection{Novelty Score $N$: Regime Coverage 制度覆盖度}

\begin{definition}[Macroeconomic Novelty Score 宏观经济新颖性评分]
\label{def:N_macro}
The novelty score $N(f_m)$ quantifies the fraction of economic regimes for which the model has been validated:
\begin{equation}
N(f_m) = \frac{|\mathcal{R}_{\mathrm{validated}}(f_m) \setminus \mathcal{R}_{\mathrm{standard}}|}{|\mathcal{R}_{\mathrm{target}}|},
\label{eq:N_macro}
\end{equation}
where:
\begin{itemize}[nosep]
    \item $\mathcal{R}_{\mathrm{validated}}(f_m) \subseteq \mathcal{R}$: the set of regimes for which the model's forecasting performance has been evaluated and published;
    \item $\mathcal{R}_{\mathrm{standard}} = \{r_{\mathrm{expansion}}\}$: the baseline regime that all forecasting models are expected to handle (post-WWII expansionary periods);
    \item $\mathcal{R}_{\mathrm{target}} = \mathcal{R} \setminus \mathcal{R}_{\mathrm{standard}}$: the set of non-standard regimes whose coverage adds novelty credit.
\end{itemize}

A model validated only on expansionary periods scores $N = 0$ (no novelty). A model validated on recessions, financial crises, and pandemics scores $N \to 1$.
\end{definition}

\subsection{\Cercis{} Score 综合评分}

\begin{definition}[\Cercis{} Score for Macroeconomic Forecasting]
\label{def:cercis_macro}
\begin{equation}
S(f_m) = \tilde{Q}(f_m) + \eta \cdot N(f_m),
\label{eq:S_macro}
\end{equation}
where $\eta \geq 0$ is a user-specified novelty-accuracy tradeoff parameter:
\begin{itemize}[nosep]
    \item $\eta = 0$: pure accuracy ranking (appropriate for operational forecasting where reliability on known regimes is paramount).
    \item $\eta > 0$: rewards models validated on diverse, non-standard regimes (appropriate for method development and stress testing).
    \item $\eta \gg 1$: prioritizes novel regime coverage over accuracy (appropriate for exploratory research).
\end{itemize}
\end{definition}

\begin{remark}[The Cercis Score and the Unidentifiability Theorem]
\label{rem:cercis_unident}
The \Cercis{} score embeds the lesson of Theorem~\ref{thm:macro_unident} in its design: the novelty term $N$ explicitly rewards models that have been tested in regimes where structural breaks are known to have occurred (recessions, crises, pandemics). A model with $N = 0$ (no non-standard regime validation) cannot credibly claim robustness to structural breaks, and its post-break errors are fully subject to the unidentifiability theorem---they could be model failure or structural change, with no validation data to distinguish them.
\end{remark}

\subsection{Illustrative \Cercis{} Rankings 示例性Cercis排名}

\begin{table}[htbp]
\centering
\caption{Illustrative \Cercis{} scores for macroeconomic forecasting models. $\tilde{Q}$ based on FRED data, 1985--2023. $\eta = 0.5$ for illustration.}
\label{tab:cercis_macro}
\begin{tabular}{@{}lcccccc@{}}
\toprule
Model / Expert & $\tilde{Q}_{\mathrm{RMSE}}$ & $\tilde{Q}_{\mathrm{cal}}$ & $\tilde{Q}$ (combined) & $N$ & $S$ \\
\midrule
Smets-Wouters DSGE~\cite{SmetsWouters2007}        & 0.62 & 0.45 & 0.56 & 0.15 & 0.64 \\
Medium-scale BVAR (Minnesota prior)~\cite{Banbura2010} & 0.71 & 0.52 & 0.65 & 0.20 & 0.75 \\
TVP-VAR~\cite{Primiceri2005}                       & 0.68 & 0.55 & 0.64 & 0.30 & 0.79 \\
ABM (Dosi et al.)~\cite{Dosi2010}                   & 0.40 & 0.20 & 0.33 & 0.50 & 0.58 \\
Gradient Boosting (FRED-MD)~\cite{Medeiros2021}     & 0.78 & 0.60 & 0.72 & 0.10 & 0.77 \\
LSTM Neural Network~\cite{Coulombe2020}             & 0.74 & 0.48 & 0.65 & 0.15 & 0.73 \\
Random Walk (benchmark)                             & 0.45 & 0.30 & 0.40 & 0.00 & 0.40 \\
\midrule
\textbf{\Yajie{} Consensus} (DSGE+VAR+ABM+ML)       & \textbf{0.82} & \textbf{0.68} & \textbf{0.78} & \textbf{0.50} & \textbf{1.03} \\
\bottomrule
\end{tabular}
\end{table}

\noindent The \Yajie{} consensus (detailed in \S6) achieves the highest \Cercis{} score by combining complementary model strengths and covering all regimes. Note that no individual model achieves $S > 0.80$; the multi-model consensus is necessary for high-quality macroeconomic forecasting. \limitationTag{2} Scores are illustrative estimates; a definitive ranking requires recomputing all methods on the standardized benchmark suite specified in \S7.

% ======================================================================
\section{Multi-Model \Yajie{} Consensus for Macroeconomic Forecasting 宏观经济多模型Yajie共识}
\label{sec:yajie}
% ======================================================================

\Yajie{} is the \SCX{} game-theoretic consensus mechanism. In macroeconomic forecasting, \Yajie{} operates at two levels: (i) within-family consensus among experts of the same modeling paradigm, and (ii) cross-family consensus aggregating DSGE, VAR, ABM, and ML forecasts into a single macroeconomic prediction.

\subsection{Two-Level Consensus Architecture 双层共识架构}

\begin{definition}[Macroeconomic \Yajie{} Game 宏观经济Yajie博弈]
\label{def:yajie_macro}
The macroeconomic \Yajie{} game is a hierarchical game with two levels:

\textbf{Level 1 --- Within-Family Consensus (模型族内共识):}
\begin{itemize}
    \item \textbf{Players:} $M_k$ experts within model family $k \in \{\mathrm{DSGE}, \mathrm{VAR}, \mathrm{ABM}, \mathrm{ML}\}$.
    \item \textbf{States:} Economic state $\statevec_t$ at each evaluation time $t$.
    \item \textbf{Actions:} Expert $m$ in family $k$ outputs a trust score $\tau_m^{(k)}(\statevec_t) \in [0, 1]$ for its own forecast.
    \item \textbf{Payoff:} $\pi_m^{(k)} = \frac{1}{T} \sum_t \ind{\tau_m^{(k)}(\statevec_t) \geq \overline{\mathrm{med}}(\boldsymbol{\tau}_{-m}^{(k)}(\statevec_t))} \cdot \mathrm{acc}_m^{(k)}(\statevec_t)$, where $\mathrm{acc}_m^{(k)}$ is the realized forecast accuracy (negative RMSE). Experts are rewarded for accurate forecasts and penalized for overconfident trust scores.
\end{itemize}

\textbf{Level 2 --- Cross-Family Consensus (跨模型族共识):}
\begin{itemize}
    \item \textbf{Players:} Four family-level aggregated forecasts $\bar{\mathbf{y}}^{(k)}$, each weighted by within-family trust scores:
    \begin{equation}
    \bar{\mathbf{y}}_t^{(k)} = \frac{\sum_{m=1}^{M_k} \tau_m^{(k)}(\statevec_t) \cdot \hat{\mathbf{y}}_t^{(m,k)}}{\sum_{m=1}^{M_k} \tau_m^{(k)}(\statevec_t)}.
    \label{eq:family_aggregate}
    \end{equation}
    \item \textbf{States:} Cross-family consensus states at each $t$.
    \item \textbf{Actions:} Family $k$ outputs a \textbf{structural confidence score} $\sigma_k(\statevec_t) \in [0, 1]$ indicating how well the current economic state matches the family's structural assumptions.
    \item \textbf{Payoff:} Families are rewarded when their structural confidence predicts post-hoc forecast accuracy. Specifically:
    \begin{equation}
    \pi_{\mathrm{cross}}^{(k)} = \Corr_t\!\left(\sigma_k(\statevec_t), -\RMSE_t^{(k)}\right),
    \label{eq:cross_payoff}
    \end{equation}
    the correlation between structural confidence and realized accuracy (higher correlation means the family knows when it will be accurate).
\end{itemize}
\end{definition}

\subsection{Consensus Mechanism 共识机制}

\begin{algorithm}[htbp]
\caption{Macroeconomic \Yajie{} Consensus with Structural Break Detection}
\label{alg:yajie_macro}
\begin{algorithmic}[1]
\Require Economic state $\statevec_t$, experts $\{f_m^{(k)}\}_{k \in \cK, m=1}^{M_k}$, historical forecast errors
\Ensure Consensus forecast $\bar{\mathbf{y}}_t$, anomaly flags, structural break probability $p_{\mathrm{break}}^{(t)}$

\State \textbf{Stage 1: Within-family aggregation 模型族内聚合}
\For{each family $k \in \{\mathrm{DSGE}, \mathrm{VAR}, \mathrm{ABM}, \mathrm{ML}\}$}
    \For{each expert $m = 1, \ldots, M_k$}
        \State Compute forecast $\hat{\mathbf{y}}_t^{(m,k)} \gets f_m^{(k)}(\statevec_t)$
        \State Compute trust score $\tau_m^{(k)} \gets \text{calibrated\_confidence}(\statevec_t, f_m^{(k)})$
    \EndFor
    \State Compute family consensus $\bar{\mathbf{y}}_t^{(k)} \gets \text{weighted\_mean}(\{\hat{\mathbf{y}}_t^{(m,k)}\}, \{\tau_m^{(k)}\})$
    \State Flag within-family anomalies: if $\max_{i,j} d_{ij}^{(t)} > \tau_{\mathrm{within}}$, mark family $k$
\EndFor

\State \textbf{Stage 2: Cross-family consensus 跨模型族共识}
\State Compute $4 \times 4$ cross-family discrepancy matrix $\mathbf{D}^{(t)}$ with entries $D_{kl}^{(t)} = \norm{\bar{\mathbf{y}}_t^{(k)} - \bar{\mathbf{y}}_t^{(l)}}_{\Sigma}$
\State Compute structural confidence $\sigma_k^{(t)}$ for each family $k$
\State \textbf{Cross-family consensus forecast:}
\begin{equation}
\bar{\mathbf{y}}_t^{\Yajie{}} = \frac{\sum_{k=1}^4 \sigma_k^{(t)} \cdot \bar{\mathbf{y}}_t^{(k)}}{\sum_{k=1}^4 \sigma_k^{(t)}}
\label{eq:yajie_consensus}
\end{equation}

\State \textbf{Stage 3: Structural break signal 结构突变信号}
\State Compute cross-family dispersion: $\delta_t = \frac{1}{6} \sum_{k < l} D_{kl}^{(t)}$
\State Compute structural break probability:
\begin{equation}
p_{\mathrm{break}}^{(t)} = \Phi\!\left( \frac{\delta_t - \mu_{\delta}}{\sigma_{\delta}} \right),
\label{eq:p_break}
\end{equation}
where $\mu_{\delta}, \sigma_{\delta}$ are the historical mean and standard deviation of cross-family dispersion, and $\Phi$ is the standard normal CDF.

\State \textbf{Theorem~\ref{thm:macro_noise} check:} For any flagged expert $m$, compute $\Pbb(\text{miss})$ bound. If $\Pbb(\text{miss}) < \alpha$, mark as ``confidently anomalous.''

\State \textbf{Theorem~\ref{thm:macro_unident} annotation:} If $p_{\mathrm{break}}^{(t)} > 0.5$ but within-family consensus holds, annotate: ``High cross-family dispersion with within-family consensus is consistent with a structural break, but model misspecification cannot be ruled out without declared assumptions (Theorem~2).''

\State \Return Consensus forecast $\bar{\mathbf{y}}_t^{\Yajie{}}$, anomaly flags, $p_{\mathrm{break}}^{(t)}$, assumption annotations
\end{algorithmic}
\end{algorithm}

\subsection{Why Cross-Family Dispersion Signals Structural Breaks 为什么跨族分歧信号化了结构突变}

\begin{proposition}[Cross-Family Dispersion as a Structural Break Indicator 跨族分歧作为结构突变指标]
\label{prop:dispersion_break}
Under \assumptionTag{3} (diverse structural assumptions across families), the expected cross-family dispersion $\E[\delta_t]$ increases at a structural break:
\begin{equation}
\E[\delta_T \mid \text{break}] > \E[\delta_{T-1} \mid \text{no break}],
\label{eq:dispersion_increase}
\end{equation}
because different model families extrapolate different aspects of pre-break dynamics, and these extrapolations diverge when the DGP changes.
\end{proposition}

\begin{proof}\rigorSketch
Let $\regime_{T-1} = r$, $\regime_T = r' \neq r$. For family $k$, the pre-break forecast $\bar{\mathbf{y}}_{T-1}^{(k)} = \E_k[\mathbf{y}_T \mid \statevec_{T-1}, \regime = r]$ is a function of the family's structural model of regime $r$. Since families make different structural assumptions about $r$, their forecasts of $\mathbf{y}_T$ differ even under regime continuity. At the break, the true DGP switches to $r'$, and each family's forecast is $\E_k[\mathbf{y}_T \mid \statevec_{T-1}, \regime = r]$ (they cannot condition on $r'$ because it is unprecedented). The cross-family dispersion $\delta_T$ is the average distance among these counterfactual forecasts. Since families differ in how regime $r$ maps to forecasts, $\delta_T$ represents the ``disagreement about what would have happened without the break''---which is mechanically larger when the break causes realized outcomes to diverge from all pre-break model predictions. $\square$
\end{proof}

\begin{remark}[COVID-19 Cross-Family Dispersion in Practice COVID-19跨族分歧的实践]
\label{rem:covid_dispersion}
During Q1--Q2 2020, DSGE models (which lack epidemiological dynamics) forecasted mild GDP declines based on historical recession patterns; VAR models extrapolated the sharp drop in high-frequency indicators (e.g., weekly unemployment claims) and forecasted severe contractions; ML models, depending on training data inclusion of the 2008 crisis, produced intermediate forecasts. The cross-family dispersion $\delta_t$ spiked to 3--4 standard deviations above its historical mean, correctly signaling a structural break even before realized GDP data confirmed the severity of the contraction.
\end{remark}

\subsection{The Consensus Forecast and Its Properties 共识预测及其性质}

\begin{proposition}[Consensus Accuracy Under Regime Stability 制度稳定下的共识准确度]
\label{prop:consensus_accuracy}
When the economic regime is stable ($\regime_t = \regime_{t-1}$), the \Yajie{} consensus forecast is asymptotically at least as accurate as the best individual family forecast:
\begin{equation}
\lim_{T \to \infty} \RMSE(\bar{\mathbf{y}}^{\Yajie{}}) \leq \min_{k \in \{1,\ldots,4\}} \RMSE(\bar{\mathbf{y}}^{(k)}) + O(T^{-1/2}).
\label{eq:consensus_bound}
\end{equation}
\end{proposition}

\begin{proof}\rigorPartial
This follows from standard forecast combination results~\cite{BatesGranger1969, Timmermann2006}. Under regime stability, the structural confidence weights $\sigma_k$ converge to weights proportional to inverse MSE. The weighted average achieves MSE no worse than the best individual forecast (up to estimation error in the weights). The $O(T^{-1/2})$ term reflects the convergence rate of the weight estimates. $\square$
\end{proof}

% ======================================================================
\section{Experimental Protocol 实验协议}
\label{sec:experiment}
% ======================================================================

We specify an experimental protocol for evaluating SCX-audited macroeconomic forecasting, designed for reproducibility using publicly available FRED data.

\subsection{Data: FRED-MD and FRED-QD 数据源}

\begin{table}[htbp]
\centering
\caption{Target macroeconomic variables and data sources. All series available from FRED (Federal Reserve Economic Data).}
\label{tab:data}
\begin{tabular}{@{}llll@{}}
\toprule
Variable & Symbol & FRED Code & Frequency \\
\midrule
Real GDP Growth            & $\GDP$  & GDPC1 (\% change)  & Quarterly \\
CPI Inflation              & $\CPI$  & CPIAUCSL (\% change) & Monthly $\to$ Quarterly \\
Unemployment Rate          & $\UNEMP$ & UNRATE          & Monthly $\to$ Quarterly \\
Federal Funds Rate         & $\FFR$  & FEDFUNDS         & Monthly $\to$ Quarterly \\
\midrule
\multicolumn{4}{@{}l}{\textbf{Conditioning variables (FRED-MD, 128 series):}} \\
Industrial Production      & IP     & INDPRO           & Monthly \\
Real Personal Consumption  & PCE    & PCEC96           & Monthly \\
S\&P 500 Index             & SP500  & SP500            & Monthly \\
10Y-3M Treasury Spread     & TS     & T10Y3M           & Monthly \\
Financial Conditions Index & NFCI   & NFCI             & Weekly $\to$ Monthly \\
\bottomrule
\end{tabular}
\end{table}

\subsection{Evaluation Periods and Structural Break Test Cases 评估期与结构突变测试}

We define three evaluation periods, each centered on a structural break:

\begin{enumerate}[label=\textbf{EP\arabic*:}, leftmargin=*]
    \item \textbf{Pre-GFC Expansion (1995Q1--2007Q4):} Baseline regime with the ``Great Moderation''---low, stable inflation and moderate GDP growth. Training data: 1960Q1--1994Q4.
    
    \item \textbf{Global Financial Crisis (2008Q1--2010Q4):} Structural break test case 1. The collapse of Lehman Brothers (September 2008) triggered the deepest recession since the Great Depression. Key break point: $T_{\mathrm{GFC}} = 2008\mathrm{Q}3$. Training data: 1960Q1--2007Q2 (excludes crisis).
    
    \item \textbf{COVID-19 Pandemic (2020Q1--2022Q4):} Structural break test case 2. The global pandemic triggered an unprecedented economic shutdown and subsequent reopening with supply chain disruptions and inflation surge. Key break points: $T_{\mathrm{COVID}}^{(1)} = 2020\mathrm{Q}1$ (lockdown), $T_{\mathrm{COVID}}^{(2)} = 2021\mathrm{Q}2$ (inflation surge). Training data: 1960Q1--2019Q4 (excludes pandemic).
\end{enumerate}

\subsection{Forecasting Models (Expert Ensemble) 预测模型（专家集合）}

\begin{enumerate}[label=\textbf{E\arabic*:}, leftmargin=*]
    \item \textbf{E1:} Smets-Wouters (2007) medium-scale DSGE with 7 observables, estimated via Bayesian methods~\cite{SmetsWouters2007}.
    \item \textbf{E2:} Christiano-Eichenbaum-Evans (2005) DSGE with financial accelerator~\cite{ChristianoEichenbaumEvans2005}.
    \item \textbf{E3:} Small-scale New Keynesian DSGE (3-equation model: IS curve, Phillips curve, Taylor rule).
    \item \textbf{E4:} BVAR with Minnesota prior, 4 variables + 128 exogenous predictors~\cite{Banbura2010}.
    \item \textbf{E5:} TVP-VAR with stochastic volatility, 4 variables~\cite{Primiceri2005}.
    \item \textbf{E6:} Factor-Augmented VAR (FAVAR) with 3 factors extracted from FRED-MD.
    \item \textbf{E7:} Large BVAR with 20 variables and adaptive hierarchical shrinkage.
    \item \textbf{E8:} Agent-Based Model: Dosi et al.\ (2010) ``Schumpeter meeting Keynes'' model~\cite{Dosi2010}.
    \item \textbf{E9:} Agent-Based Model: Mark-0 model with banking sector~\cite{HaldaneTurrell2018}.
    \item \textbf{E10:} Gradient Boosting (XGBoost) on FRED-MD with 128 predictors~\cite{Medeiros2021}.
    \item \textbf{E11:} LSTM recurrent neural network with attention mechanism~\cite{Coulombe2020}.
    \item \textbf{E12:} Random Forest on FRED-MD with macroeconomic factor features.
    \item \textbf{E13:} Transformer-based time series model (Informer architecture).
    \item \textbf{E14:} Professional forecaster consensus (Survey of Professional Forecasters, SPF)---serves as a human-expert baseline and ground-truth anchor.
\end{enumerate}

Total: $M_{\mathrm{DSGE}} = 3$, $M_{\mathrm{VAR}} = 4$, $M_{\mathrm{ABM}} = 2$, $M_{\mathrm{ML}} = 4$, plus 1 anchor. $M_{\mathrm{cross}}^{(i)} \geq 9$ for all experts, providing reasonable detection power per Corollary~\ref{cor:required_M}.

\subsection{Evaluation Metrics 评估指标}

\begin{enumerate}[label=(\roman*)]
    \item \textbf{RMSE:} Root Mean Squared Error, computed per-variable and averaged:
    \begin{equation}
    \RMSE^{(m)} = \frac{1}{T H d_y} \sum_{t=1}^T \sum_{h=1}^H \norm{\hat{y}_{t+h}^{(m)} - y_{t+h}}_{\Sigma}.
    \end{equation}
    
    \item \textbf{CRPS:} Continuous Ranked Probability Score for probabilistic forecasts (Eq.~\ref{eq:CRPS}).
    
    \item \textbf{Cross-family dispersion $\delta_t$:} Average pairwise discrepancy across model families (Eq.~\ref{eq:dispersion_increase} context).
    
    \item \textbf{Structural break detection accuracy:} For known break points $T_{\mathrm{GFC}}$ and $T_{\mathrm{COVID}}$, compute $p_{\mathrm{break}}^{(t)}$ from Eq.~\ref{eq:p_break}. Report AUC (area under the ROC curve) for binary classification of break vs.\ non-break periods.
    
    \item \textbf{Consensus forecast error:} $\RMSE(\bar{\mathbf{y}}^{\Yajie{}})$ compared against best individual expert and simple equal-weighted average.
    
    \item \textbf{Assumption audit completeness:} For each error attribution claim, verify that at least one assumption from Corollary~\ref{cor:assumption_mandate_macro} is explicitly declared. Score: fraction of claims with complete assumption declarations.
\end{enumerate}

\subsection{Experimental Procedure 实验流程}

\begin{algorithm}[htbp]
\caption{SCX Macroeconomic Forecast Audit Protocol}
\label{alg:macro_audit}
\begin{algorithmic}[1]
\Require FRED data $\cD_{\mathrm{all}}$, experts $\{f_m\}_{m=1}^{14}$, evaluation periods EP1--EP3, break points $T_{\mathrm{GFC}}, T_{\mathrm{COVID}}^{(1)}, T_{\mathrm{COVID}}^{(2)}$
\Ensure Audit report with $\{Q(f_m), N(f_m), S(f_m)\}$, Yajie consensus, structural break signals, assumption declarations

\For{each evaluation period EP $p \in \{1, 2, 3\}$}
    \State Split data: $\cD_{\mathrm{train}}^{(p)} \gets$ pre-period data, $\cD_{\mathrm{test}}^{(p)} \gets$ period data
    \State \textbf{Train experts:} For each expert $m$, estimate/calibrate on $\cD_{\mathrm{train}}^{(p)}$
    
    \For{each time $t \in \cD_{\mathrm{test}}^{(p)}$}
        \State Construct state $\statevec_t$ from data up to $t$
        \For{each expert $m = 1, \ldots, M$}
            \State Compute forecast $\hat{\mathbf{y}}_{t+1:t+H}^{(m)} \gets f_m(\statevec_t)$
        \EndFor
        
        \State \textbf{Yajie Level 1:} Within-family consensus, compute $\bar{\mathbf{y}}_t^{(k)}$, flag anomalies
        \State \textbf{Yajie Level 2:} Cross-family consensus, compute $\bar{\mathbf{y}}_t^{\Yajie{}}$
        \State \textbf{Structural break signal:} Compute $\delta_t$, $p_{\mathrm{break}}^{(t)}$
        
        \State \textbf{Theorem~\ref{thm:macro_noise} check:} For each expert, compute $\Pbb(\text{miss})$ bound
        \State \textbf{Theorem~\ref{thm:macro_unident} annotation:} If $p_{\mathrm{break}}^{(t)} > 0.5$, annotate with assumption requirement
    \EndFor
    
    \State Compute per-expert $\RMSE^{(m)}$, $\CRPS^{(m)}$, $Q(f_m)$, $N(f_m)$, $S(f_m)$
    \State Compute consensus $\RMSE$, structural break detection AUC
    \State \textbf{Assumption audit:} For each period with $p_{\mathrm{break}}^{(t)} > 0.5$, verify assumption declarations
\EndFor

\State \Return Comprehensive audit report
\end{algorithmic}
\end{algorithm}

\subsection{Expected Results and Hypotheses 预期结果与假设}

Based on the theoretical framework, we formulate the following testable hypotheses:

\begin{enumerate}[label=\textbf{H\arabic*:}, leftmargin=*]
    \item \textbf{Cross-family dispersion increases at structural breaks:} $\delta_t$ at $t \in \{2008\mathrm{Q}3, 2020\mathrm{Q}1\}$ exceeds the 95th percentile of its historical distribution (pre-2008).
    
    \item \textbf{Yajie consensus outperforms individual families during structural breaks:} $\RMSE(\bar{\mathbf{y}}^{\Yajie{}}) < \min_k \RMSE(\bar{\mathbf{y}}^{(k)})$ at $T_{\mathrm{GFC}}$ and $T_{\mathrm{COVID}}$.
    
    \item \textbf{DSGE models have the largest post-break errors:} Due to their strong structural assumptions, DSGE forecasts exhibit the largest $\varepsilon_{\mathrm{break}} + \varepsilon_{\mathrm{spec}}$ during unprecedented regime changes, while ML and VAR models adapt more flexibly (at the cost of less structural interpretability).
    
    \item \textbf{Within-family consensus masks structural vulnerability:} During the pre-GFC expansion (EP1), within-family consensus is high for all families, but cross-family dispersion is non-zero, reflecting fundamentally different structural assumptions. This cross-family disagreement is the early warning signal that within-family consensus cannot provide.
    
    \item \textbf{Theorem~\ref{thm:macro_unident} is empirically binding:} For both GFC and COVID break points, no purely data-driven criterion can assign the post-break forecast error uniquely to model failure vs.\ structural change. Only models with explicit, pre-registered assumption declarations produce falsifiable error attributions.
\end{enumerate}

% ======================================================================
\section{Discussion 讨论}
\label{sec:discussion}
% ======================================================================

\subsection{Implications for Central Bank Forecasting 对央行预测的启示}

Central banks (Federal Reserve, ECB, Bank of England, People's Bank of China 中国人民银行) produce macroeconomic forecasts that directly inform monetary policy decisions. The \SCX{} framework has three immediate implications for central bank forecasting practice:

\begin{enumerate}[label=(\roman*)]
    \item \textbf{Multi-model ensembles are not optional.} Theorem~\ref{thm:macro_noise} proves that single-model forecasting (or multi-model forecasting using only one model family) provides no formal guarantee against missing systematic errors. Central banks already use multiple models (FRB/US, SIGMA, EDO at the Fed; NAWM, New Area-Wide Model at the ECB), but these are predominantly DSGE variants. The detection guarantee requires cross-family diversity: at minimum, DSGE + VAR + ML.
    
    \item \textbf{Structural break signals should be systematically monitored.} The cross-family dispersion $\delta_t$ (Proposition~\ref{prop:dispersion_break}) provides a real-time, model-based indicator of potential structural breaks. Central banks monitor financial conditions indices and uncertainty measures; we propose adding $\delta_t$ as a \textit{model disagreement index} (模型分歧指数) to the dashboard of leading indicators.
    
    \item \textbf{Honest uncertainty communication requires assumption declarations.} When a central bank's forecast deviates from realized outcomes, the post-hoc explanation---``the economy changed in ways our models didn't anticipate''---is scientifically incomplete without declaring the assumption under which this attribution is made. Theorem~\ref{thm:macro_unident} provides the formal framework for this declaration.
\end{enumerate}

\subsection{The Honest Forecaster Principle 诚实预测者原则}

\begin{proposition}[The Honest Macroeconomic Forecaster 诚实的宏观经济预测者]
\label{prop:honest_forecaster}
An \textbf{honest macroeconomic forecast} is a triple:
\begin{equation}
\mathcal{H}_t = \left( \bar{\mathbf{y}}_t^{\Yajie{}}, \; \mathcal{A}_t, \; \mathcal{U}_t \right),
\label{eq:honest_forecast}
\end{equation}
where:
\begin{itemize}[nosep]
    \item $\bar{\mathbf{y}}_t^{\Yajie{}}$ is the multi-model consensus point forecast (and its predictive distribution);
    \item $\mathcal{A}_t$ is the set of \textbf{declared assumptions} under which the forecast is valid (model scope, regime continuity, measurement error bounds);
    \item $\mathcal{U}_t$ is the set of \textbf{acknowledged unidentifiabilities}---claims about the forecast error that cannot be verified without further assumptions (per Theorem~\ref{thm:macro_unident}).
\end{itemize}
The forecast is \textbf{honest} if and only if $\mathcal{A}_t$ and $\mathcal{U}_t$ are non-empty and explicitly stated.
\end{proposition}

This principle operationalizes the SCX ``Honest Agent'' theorem (\ThmSCXHonest) for macroeconomic forecasting. A forecast without declared assumptions and acknowledged unidentifiabilities is not necessarily wrong, but it is \textbf{scientifically incomplete}---it makes implicit claims that cannot be evaluated.

\subsection{Relationship to Existing Forecasting Methods 与现有预测方法的关系}

\textbf{Forecast combination \cite{BatesGranger1969, Timmermann2006}.} The extensive literature on combining forecasts establishes that weighted averages of diverse models outperform individual models. The \SCX{} framework provides the missing theoretical justification for \textit{why} diversity matters: it is not merely that diverse models have uncorrelated errors (the standard argument), but that cross-family diversity is \textit{necessary} for detecting systematic errors (Theorem~\ref{thm:macro_noise}) and for signaling structural breaks (Proposition~\ref{prop:dispersion_break}).

\textbf{Bayesian model averaging (BMA) \cite{ClarkMcCracken2010}.} BMA weights forecasts by posterior model probabilities, providing a coherent probabilistic framework. However, BMA assumes the true model is in the model set---an assumption violated at structural breaks. The \Yajie{} consensus does not require this assumption; it explicitly models the possibility that \textit{all} models are misspecified post-break and uses cross-family dispersion as the detection signal.

\textbf{Real-time monitoring of structural change \cite{GiacominiRossi2009}.} Giacomini and Rossi (2009) and subsequent work develop tests for forecast breakdown. The \SCX{} framework complements these tests with a theoretical guarantee on detection probability (Theorem~\ref{thm:macro_noise}) and a formal statement of the residual unidentifiability (Theorem~\ref{thm:macro_unident}).

\textbf{Scenario-based forecasting \cite{Wieland2020}.} Central banks increasingly use scenario analysis---alternative paths conditional on different assumptions about the pandemic, fiscal policy, or energy prices. The \SCX{} framework formalizes this as assumption-declared forecasting: each scenario is a forecast $\bar{\mathbf{y}}_t^{\Yajie{}}$ under a specific declared assumption set $\mathcal{A}_t$, and the set of scenarios spans the space of acknowledged unidentifiabilities $\mathcal{U}_t$.

\subsection{Limitations 局限性}

We explicitly acknowledge the following limitations:

\begin{enumerate}[label=\limitationTag{\arabic*}, leftmargin=*, resume]
    \item \textbf{(Computational Cost)} Fielding 14 forecasting experts (DSGE estimation, BVAR with large datasets, ABM simulation, ML training) is computationally demanding. For a central bank with dedicated research staff, this is feasible; for an individual researcher, it may be prohibitive. We encourage the development of pre-trained, publicly available expert ensembles.
    
    \item \textbf{(Regime Classification Lag)} \assumptionTag{1} acknowledges that regimes are identifiable only with delay. The real-time structural break signal $p_{\mathrm{break}}^{(t)}$ (Eq.~\ref{eq:p_break}) provides an instantaneous indicator, but definitive regime assignment requires historical perspective.
    
    \item \textbf{(Within-Family Expert Diversity)} The DSGE and ABM families have fewer ``off-the-shelf'' variants than VAR (many prior specifications) or ML (many architectures). Ensuring genuine within-family diversity (Assumption~\ref{ass:A5} in the CFD paper analog) requires careful model construction.
    
    \item \textbf{(Measurement Error in Real-Time Data)} Macroeconomic data are revised substantially after initial release. GDP growth for a given quarter may be revised by 1--2 percentage points between the advance estimate and the final revision years later. This injects additional unidentifiability (realized outcomes differ from true economic states) that our framework inherits.
    
    \item \textbf{(Cercis Score Subjectivity)} The regime weights $w_r$ and novelty-accuracy tradeoff $\eta$ embed normative judgments about which regimes matter most. For inflation-targeting central banks, $w_{\mathrm{stagflation}}$ should be high; for employment-focused mandates, $w_{\mathrm{recession}}$ should be high. The framework is flexible but requires explicit justification of weight choices.
    
    \item \textbf{(Lucas Critique Recursion)} The Lucas Critique applies to the \SCX{} framework itself: if forecasters know they are being audited by \SCX{}, they may change their behavior (strategic forecast reporting). This is a meta-level unidentifiability that Theorem~\ref{thm:macro_unident} does not resolve---it is the subject of future work on ``auditing the auditors.''
\end{enumerate}

\subsection{Future Work 未来工作}

\begin{enumerate}[label=(\roman*)]
    \item \textbf{Empirical validation:} Implement the experimental protocol (\S7) on FRED data with the 14-expert ensemble. Measure detection probabilities, structural break signal accuracy, and consensus forecast performance.
    
    \item \textbf{Real-time structural break monitoring system:} Deploy the cross-family dispersion $\delta_t$ as a continuously updated public indicator, analogous to the St.\ Louis Fed's Financial Stress Index.
    
    \item \textbf{Chinese macroeconomic forecasting 中国宏观经济预测:} Adapt the framework to China's institutional context, where the policy regime (五年规划 Five-Year Plans, 人民银行 PBoC policy tools) and data structure (National Bureau of Statistics 国家统计局 releases) differ from Western economies. This requires defining China-specific regimes $\mathcal{R}_{\mathrm{CN}}$ and model families suitable for the Chinese economy.
    
    \item \textbf{Assumption declaration standards:} Develop standardized templates for macroeconomic forecast assumption declarations (analogous to the STROBE statement in epidemiology or the CONSORT statement in clinical trials), enabling systematic auditing of forecast honesty.
    
    \item \textbf{Generalization to international macroeconomics:} Extend the framework to multi-country forecasting with spillover effects, where structural breaks in one country (e.g., Brexit, China's COVID-zero policy) propagate to others through trade and financial linkages.
\end{enumerate}

% ======================================================================
\section{Conclusion 结论}
\label{sec:conclusion}
% ======================================================================

We have applied the \SCX{} quality audit framework to macroeconomic forecasting, providing formal guarantees for forecast error detection, a fundamental unidentifiability result linking the Lucas Critique to all structural breaks, and a practical consensus mechanism for combining DSGE, VAR, ABM, and ML forecasts.

Three theorems anchor the contribution:

\begin{enumerate}[label=(\roman*)]
    \item \textbf{Theorem~\ref{thm:macro_noise} (Multi-Model Forecast Error Detection):} The probability of missing a systematic forecasting error of magnitude $\Delta$ decays exponentially in the number of cross-family experts: $\Pbb(\text{miss}) \leq \exp(-2 M_{\mathrm{cross}} (\gamma - \theta)^2)$. This provides the first formal guarantee for multi-model macroeconomic forecasting.
    
    \item \textbf{Theorem~\ref{thm:macro_unident} (Model Misspecification vs.\ Structural Break Unidentifiability):} When a forecasting model's predictions deviate from realized outcomes, the attribution of this error to model failure versus genuine structural change is logically underdetermined without declared assumptions. The Lucas Critique emerges as the special case where the structural break is an observed policy regime change.
    
    \item \textbf{Proposition~\ref{prop:dispersion_break} (Cross-Family Dispersion as Structural Break Signal):} Cross-family forecast dispersion increases at structural breaks because different model families extrapolate different aspects of pre-break dynamics, providing a model-based, real-time structural break indicator.
\end{enumerate}

The \Cercis{} score $S = \tilde{Q} + \eta N$ operationalizes the quality-novelty tradeoff, and the \Yajie{} two-level consensus mechanism (within-family, then cross-family) provides a practical procedure for multi-model forecasting with explicit assumption tracking.

\textbf{The central message is methodological:} Macroeconomic forecasting under structural breaks requires not only better models, but a better framework for acknowledging when models cannot distinguish their own failure from a changing world. The \SCX{} framework provides exactly this: a formal structure for declaring assumptions, detecting errors, and honestly communicating the limits of macroeconomic knowledge.

\medskip
\noindent\textbf{核心命题 核心命题 (Core Proposition in Chinese):}

宏观经济预测面临的根本挑战不是模型精度不足，而是结构突变（金融危机、疫情、政策转向）使得任何基于历史数据训练的模型在突变点系统性失效。SCX审计框架通过三个机制应对这一挑战：(1) 多模型家族共识检测——当DSGE、VAR、基于主体的模型和机器学习模型同时部署时，遗漏系统性预测误差的概率随跨族专家数量指数衰减；(2) 不可辨识性定理——证明模型误设与结构突变在观测上是等价的，卢卡斯批判是政策制度突变可被观测时的特例；(3) \Cercis{}评分——结合预测准确度（RMSE和密度校准CRPS）与制度新颖性，对预测模型进行综合排名。核心洞察：跨模型族分歧不仅是鲁棒性的最佳实践，更是检测结构突变和诚实承认预测局限性的逻辑必然。

\medskip
\noindent\textbf{关键术语中英对照 (Key Terminology Chinese-English Glossary):}

\begin{itemize}[nosep]
    \item 宏观经济预测 (Macroeconomic Forecasting)
    \item 结构突变 (Structural Break)
    \item 卢卡斯批判 (Lucas Critique)
    \item 动态随机一般均衡 (DSGE: Dynamic Stochastic General Equilibrium)
    \item 向量自回归 (VAR: Vector Autoregression)
    \item 基于主体的模型 (ABM: Agent-Based Model)
    \item 机器学习 (ML: Machine Learning)
    \item 多模型共识 (Multi-Model Consensus)
    \item 不可辨识性 (Unidentifiability)
    \item 预测误差检测 (Forecast Error Detection)
    \item 密度校准 (Density Calibration)
    \item 制度新颖性 (Regime Novelty)
    \item 跨模型族分歧 (Cross-Family Dispersion)
    \item 诚实预测者原则 (Honest Forecaster Principle)
\end{itemize}

% ======================================================================
% BIBLIOGRAPHY
% ======================================================================
\begin{thebibliography}{99}

\bibitem{Lucas1976}
R.~E.~Lucas.
\newblock Econometric policy evaluation: A critique.
\newblock {\em Carnegie-Rochester Conference Series on Public Policy}, 1:19--46, 1976.

\bibitem{Sims1980}
C.~A.~Sims.
\newblock Macroeconomics and reality.
\newblock {\em Econometrica}, 48(1):1--48, 1980.

\bibitem{SmetsWouters2007}
F.~Smets and R.~Wouters.
\newblock Shocks and frictions in US business cycles: A Bayesian DSGE approach.
\newblock {\em American Economic Review}, 97(3):586--606, 2007.

\bibitem{ChristianoEichenbaumEvans2005}
L.~J.~Christiano, M.~Eichenbaum, and C.~L.~Evans.
\newblock Nominal rigidities and the dynamic effects of a shock to monetary policy.
\newblock {\em Journal of Political Economy}, 113(1):1--45, 2005.

\bibitem{Primiceri2005}
G.~E.~Primiceri.
\newblock Time varying structural vector autoregressions and monetary policy.
\newblock {\em Review of Economic Studies}, 72(3):821--852, 2005.

\bibitem{Banbura2010}
M.~Banbura, D.~Giannone, and L.~Reichlin.
\newblock Large Bayesian vector auto regressions.
\newblock {\em Journal of Applied Econometrics}, 25(1):71--92, 2010.

\bibitem{Dosi2010}
G.~Dosi, G.~Fagiolo, and A.~Roventini.
\newblock Schumpeter meeting Keynes: A policy-friendly model of endogenous growth and business cycles.
\newblock {\em Journal of Economic Dynamics and Control}, 34(9):1748--1767, 2010.

\bibitem{HaldaneTurrell2018}
A.~G.~Haldane and A.~E.~Turrell.
\newblock An interdisciplinary model for macroeconomics.
\newblock {\em Oxford Review of Economic Policy}, 34(1-2):219--251, 2018.

\bibitem{Coulombe2020}
P.~G.~Coulombe, M.~Leroux, D.~Stevanovic, and S.~Surprenant.
\newblock How is machine learning useful for macroeconomic forecasting?
\newblock {\em Journal of Applied Econometrics}, 37(5):920--948, 2022.

\bibitem{Medeiros2021}
M.~C.~Medeiros, G.~F.~Vasconcelos, A.~Veiga, and E.~Zilberman.
\newblock Forecasting inflation in a data-rich environment: The benefits of machine learning methods.
\newblock {\em Journal of Business and Economic Statistics}, 39(1):98--119, 2021.

\bibitem{Wieland2020}
V.~Wieland and M.~Wolters.
\newblock Macroeconomic model comparisons and forecast competitions.
\newblock {\em VoxEU.org}, 2020.

\bibitem{Ng2021}
S.~Ng.
\newblock Modeling macroeconomic variations after COVID-19.
\newblock {\em NBER Working Paper}, 29060, 2021.

\bibitem{AiolfiTimmermann2006}
M.~Aiolfi and A.~Timmermann.
\newblock Persistence in forecasting performance and conditional combination strategies.
\newblock {\em Journal of Econometrics}, 135(1-2):31--53, 2006.

\bibitem{MincerZarnowitz1969}
J.~A.~Mincer and V.~Zarnowitz.
\newblock The evaluation of economic forecasts.
\newblock In {\em Economic Forecasts and Expectations}, NBER, pp.~3--46, 1969.

\bibitem{GneitingRaftery2007}
T.~Gneiting and A.~E.~Raftery.
\newblock Strictly proper scoring rules, prediction, and estimation.
\newblock {\em Journal of the American Statistical Association}, 102(477):359--378, 2007.

\bibitem{BatesGranger1969}
J.~M.~Bates and C.~W.~J.~Granger.
\newblock The combination of forecasts.
\newblock {\em Operational Research Quarterly}, 20(4):451--468, 1969.

\bibitem{Timmermann2006}
A.~Timmermann.
\newblock Forecast combinations.
\newblock {\em Handbook of Economic Forecasting}, 1:135--196, 2006.

\bibitem{ClarkMcCracken2010}
T.~E.~Clark and M.~W.~McCracken.
\newblock Averaging forecasts from VARs with uncertain instabilities.
\newblock {\em Journal of Applied Econometrics}, 25(1):5--29, 2010.

\bibitem{GiacominiRossi2009}
R.~Giacomini and B.~Rossi.
\newblock Detecting and predicting forecast breakdowns.
\newblock {\em Review of Economic Studies}, 76(2):669--705, 2009.

\bibitem{Hoeffding1963}
W.~Hoeffding.
\newblock Probability inequalities for sums of bounded random variables.
\newblock {\em Journal of the American Statistical Association}, 58(301):13--30, 1963.

\bibitem{Liang1986}
K.-Y.~Liang and S.~L.~Zeger.
\newblock Longitudinal data analysis using generalized linear models.
\newblock {\em Biometrika}, 73(1):13--22, 1986.

\bibitem{SCX2025}
SCX Research Collective.
\newblock SCX: Structured Causal eXamination Framework for Auditable AI.
\newblock {\em Technical Report}, 2025.

\bibitem{StockWatson2002}
J.~H.~Stock and M.~W.~Watson.
\newblock Macroeconomic forecasting using diffusion indexes.
\newblock {\em Journal of Business and Economic Statistics}, 20(2):147--162, 2002.

\bibitem{McCrackenNg2016}
M.~W.~McCracken and S.~Ng.
\newblock FRED-MD: A monthly database for macroeconomic research.
\newblock {\em Journal of Business and Economic Statistics}, 34(4):574--589, 2016.

\bibitem{FaustWright2013}
J.~Faust and J.~H.~Wright.
\newblock Forecasting inflation.
\newblock In {\em Handbook of Economic Forecasting}, vol.~2, pp.~2--56, Elsevier, 2013.

\bibitem{DelNegroSchorfheide2013}
M.~Del Negro and F.~Schorfheide.
\newblock DSGE model-based forecasting.
\newblock In {\em Handbook of Economic Forecasting}, vol.~2, pp.~57--140, Elsevier, 2013.

\bibitem{Hansen2016}
B.~E.~Hansen.
\newblock Efficient shrinkage in parametric models.
\newblock {\em Journal of Econometrics}, 190(1):115--132, 2016.

\bibitem{DieboldMariano1995}
F.~X.~Diebold and R.~S.~Mariano.
\newblock Comparing predictive accuracy.
\newblock {\em Journal of Business and Economic Statistics}, 13(3):253--263, 1995.

\bibitem{Cochrane2011}
J.~H.~Cochrane.
\newblock Determinacy and identification with Taylor rules.
\newblock {\em Journal of Political Economy}, 119(3):565--615, 2011.

\bibitem{NakamuraSteinsson2018}
E.~Nakamura and J.~Steinsson.
\newblock Identification in macroeconomics.
\newblock {\em Journal of Economic Perspectives}, 32(3):59--86, 2018.

\bibitem{ReinhartRogoff2009}
C.~M.~Reinhart and K.~S.~Rogoff.
\newblock {\em This Time Is Different: Eight Centuries of Financial Folly}.
\newblock Princeton University Press, 2009.

\bibitem{JordaSchularickTaylor2017}
O.~Jorda, M.~Schularick, and A.~M.~Taylor.
\newblock Macrofinancial history and the new business cycle facts.
\newblock {\em NBER Macroeconomics Annual}, 31:213--263, 2017.

\end{thebibliography}

% ======================================================================
\appendix
\section{Assumption Map 假设地图}
\label{sec:app_assumptions}
% ======================================================================

Table~\ref{tab:assumptions} provides the complete map of all assumptions declared in this paper, their types, and their dependencies.

\begin{table}[htbp]
\centering
\caption{Complete assumption map for SCX macroeconomic forecasting audit.}
\label{tab:assumptions}
\begin{tabular}{@{}llll@{}}
\toprule
Assumption & Description & Type & Depends On \\
\midrule
A0 & Motivating observation & Empirical & --- \\
A1 & Regime identifiability ex post & Structural & NBER methodology \\
A2 & Bounded macroeconomic output & Empirical & Historical bounds \\
A3 & Expert diversity across families & Structural & Model family definitions \\
A4 & Detectability condition & Structural & A3 \\
A5 & Irreducible noise unforecastable & Structural & Rational expectations \\
\midrule
L0 & Framework is descriptive, not prescriptive & Scope & --- \\
L1 & Hoeffding independence idealization & Statistical & Estimation of $\bar{\rho}$ \\
L2 & Cercis scores illustrative, not definitive & Methodological & Full benchmark recomputation \\
L3 & Regime classification lag & Structural & A1 \\
L4 & Within-family diversity constraint & Practical & Model development \\
L5 & Real-time data measurement error & Empirical & Data revision studies \\
L6 & Cercis weight subjectivity & Normative & Policy mandate \\
L7 & Lucas Critique recursion & Meta-theoretical & Future work \\
\bottomrule
\end{tabular}
\end{table}

\section{Proof of the Hoeffding Bound with Correlated Errors 相关误差下的Hoeffding界证明}
\label{sec:app_hoeffding}

We provide the derivation of the effective sample size correction used in Theorem~\ref{thm:macro_noise}.

\begin{proposition}[Effective Multiplicity Under Correlation 相关性下的有效多重性]
\label{prop:effective_M}
For $M$ experts with average pairwise error correlation $\bar{\rho}$, the variance of the mean of their error indicators satisfies:
\begin{equation}
\Var(\bar{X}) = \frac{\sigma^2}{M} \cdot (1 + (M-1)\bar{\rho}),
\label{eq:variance_inflation}
\end{equation}
where $\sigma^2 = \Var(X_j)$. Consequently, the Hoeffding bound uses effective multiplicity:
\begin{equation}
M_{\mathrm{eff}} = \frac{M}{1 + (M-1)\bar{\rho}}.
\label{eq:M_eff}
\end{equation}
\end{proposition}

\begin{proof}
For correlated Bernoulli variables $X_1, \ldots, X_M$ with common mean $\mu = \E[X_j]$ and pairwise correlation $\rho_{ij} = \Corr(X_i, X_j)$:
\begin{align}
\Var(\bar{X}) &= \frac{1}{M^2} \Var\!\left(\sum_{j=1}^M X_j\right) \\
&= \frac{1}{M^2} \left( \sum_{j=1}^M \Var(X_j) + \sum_{i \neq j} \Cov(X_i, X_j) \right) \\
&= \frac{\sigma^2}{M^2} \left( M + \sum_{i \neq j} \rho_{ij} \right) \\
&= \frac{\sigma^2}{M} \left( 1 + \frac{2}{M} \sum_{i < j} \rho_{ij} \right) \\
&= \frac{\sigma^2}{M} (1 + (M-1)\bar{\rho}),
\end{align}
where $\bar{\rho} = \frac{2}{M(M-1)} \sum_{i < j} \rho_{ij}$.

The Hoeffding inequality for the mean of correlated bounded variables can be applied with the variance-inflated sample size: the concentration is equivalent to that of $M_{\mathrm{eff}} = M / (1 + (M-1)\bar{\rho})$ independent variables. This is a standard correction in clustered and longitudinal data analysis~\cite{Liang1986}. $\square$
\end{proof}

\section{CRPS Decomposition for Density Forecast Evaluation 密度预测评估的CRPS分解}
\label{sec:app_crps}

\begin{proposition}[CRPS Decomposition 连续排名概率评分分解]
\label{prop:crps_decomp}
The CRPS can be decomposed into calibration and sharpness components:
\begin{equation}
\CRPS(\hat{P}, y) = \underbrace{\E_{\hat{P}}[|X - y|] - \frac{1}{2}\E_{\hat{P}}[|X - X'|]}_{\text{sharpness}} + \underbrace{\left( \E_{\hat{P}}[F_{\hat{P}}^{-1}(\alpha)] - F_y^{-1}(\alpha) \right)^2 \text{ term}}_{\text{calibration}},
\label{eq:crps_decomp}
\end{equation}
where $X, X' \sim \hat{P}$ are independent draws. This decomposition enables separate auditing of whether a forecasting model's predictive distribution is well-calibrated (correct coverage) and sharp (narrow intervals).
\end{proposition}

This decomposition is due to~\cite{GneitingRaftery2007} and is included for completeness in the macroeconomic context, where density forecasts (fan charts) are central to central bank communication.

\end{document}
